% Section 6, part 4.
\begin{lem}[Functions of $x$ pair to zero]
\label{lem:mod-xfun}
Let $2h-2+n>0$, $n\geq 1$, let $p_2,\dots,p_n$ be points of $\bbP^1$ that are not branch points, and let $f$ be holomorphic near the five critical values $x(\varsigma_b)$ (for instance $f\in\bbQ(v)[X]$, $v\neq 0$ fixed). Then for each $b$,
\[
\Res_{p=\varsigma_b}f(x(p))\,\omega_{h,n}(p,p_2,\dots,p_n)=0 .
\]
Consequently every pairing of $\omega_{h,n}$ in which one slot is paired with $f(x)$ by the residue pairing, the other slots being paired by residues at the branch points, on circles or along rays, vanishes; legs $\phi,\phi'\in\cL$ with $\phi-\phi'\in\bbQ(v)[x]$ have the same residue pairings; and (with $f=1$) $\omega_{h,n}$ has zero residues at the branch points in each variable.
\end{lem}

\begin{proof}
\emph{The trace identity.} For $X\in\bbP^1$ let $p^1(X),\dots,p^5(X)$ be the points of $x^{-1}(X)$, the roots of $T^5-4XT+4v$, and let $\mathrm{Tr}$ be the push-forward of $1$-forms along $x\colon\bbP^1\to\bbP^1$. We claim that $\mathrm{Tr}\,\omega(X):=\sum_{i=1}^5\omega_{h,n}(p^i(X),p_2,\dots,p_n)=0$ identically. (i) For $X$ not a critical value and $q'\notin x^{-1}(X)$, $\prod_i(q'-p^i(X))=q'^5-4Xq'+4v=4q'\,(x(q')-X)$; since $\omega_{0,2}(p,q')=\dd_p\dd_{q'}\log(p-q')$, this gives $\sum_i\omega_{0,2}(p^i(X),q')=\dd_X\dd_{q'}\log(x(q')-X)=\dd X\,\dd x(q')/(X-x(q'))^2$. (ii) Choose small discs $U_b$ around the branch points with $\iota_b(U_b)=U_b$, disjoint from each other and from $\{p_2,\dots,p_n\}$, and put $U:=\bigcup_bU_b$. For $p_1\notin U$ the recursion \eqref{eq:mod-kernel} reads $\omega_{h,n}(p_1,J)=\sum_b(2\pi\ri)^{-1}\oint_{\partial U_b}K_b(p_1,q)\,\mathcal R_b(q,J)$, $J=(p_2,\dots,p_n)$, where $\mathcal R_b$ (a local letter) is the bracket of \eqref{eq:mod-kernel}, the factors $\omega_{0,2}(q,p_l)$ included; it does not depend on $p_1$, and the integrand has no pole in $U_b$ other than $\varsigma_b$. The path from $\iota_bq$ to $q$ in $K_b$ stays in $U_b$. (iii) For $X$ outside the compact set $x(\overline U)$ all $p^i(X)$ lie outside $U$, and by (i)
\[
\sum_i\int_{\iota_bq}^q\omega_{0,2}(p^i(X),\cdot)=\dd X\Big[\frac{1}{X-x(q)}-\frac{1}{X-x(\iota_bq)}\Big]=0,
\]
because $x(\iota_bq)=x(q)$. Summing (ii) over $i$ (a finite sum, exchanged with the contour integrals) gives $\mathrm{Tr}\,\omega(X)=0$ for these $X$; as $\mathrm{Tr}\,\omega$ is a meromorphic $1$-form in $X$, it vanishes identically. (This is the computation of \cite[App.~A, (A-1), (A-5)]{Eynard:2007kz}.)

\emph{The residues.} Fix $b$ and put $X_b:=x(\varsigma_b)$. The fibre $x^{-1}(X_b)$ consists of $\varsigma_b$, of local degree $2$, and three further points $q_1,q_2,q_3$ of local degree $1$; none of them is a branch point, because the critical values are pairwise distinct, and none is $0$ or $\infty$, which are mapped to $\infty$. Hence $f(x)\,\omega_{h,n}(\cdot,J)$ is holomorphic at the $q_k$ (\cref{conv:mod-kernel}: poles only at branch points, none on the diagonals). For a $1$-form $\eta$ meromorphic near the fibre, $\Res_{X=X_b}\mathrm{Tr}(\eta)=\sum_{q\in x^{-1}(X_b)}\Res_q\eta$: the preimage of a small circle around $X_b$ is a union of small loops, one around each point of the fibre, mapped onto the circle with the local degree, so $\oint\mathrm{Tr}(\eta)=\sum_q\oint\eta$. With $\eta=f(x)\,\omega_{h,n}(\cdot,J)$ we have $\mathrm{Tr}(\eta)=f(X)\,\mathrm{Tr}\,\omega=0$, so $0=\Res_{\varsigma_b}\eta+\sum_k\Res_{q_k}\eta=\Res_{\varsigma_b}\eta$.

\emph{Consequences.} $\omega_{h,n}$ is symmetric (the paragraph after \cref{conv:mod-kernel}), so the slot carrying $f(x)$ may be taken to be the first. The contours of the other slots can be chosen disjoint from small discs around the branch points used in the first slot, and $\omega_{h,n}$ has no diagonal poles, so the residue in the first slot can be taken first, for fixed points of the other contours, and it vanishes. The other two statements follow.
\end{proof}

The statement and its deduction are also \cite[Thm.~4.4 (4-23), Cor.~4.1 (4-26)--(4-29), App.~A (A-25)]{Eynard:2007kz}; the proof above does not use them.

\subsection{The leg dictionary and the residue identity}
\label{sec:mod-legs}

\begin{defn}[The legs]
\label{def:mod-legs}
Let $\phi^b_0:=p,\ \log p,\ v/p,\ p^3,\ p^2$ for $b=0,\dots,4$. For $F$ in the span of the $p^j(\log p)^l$ ($j\in\bbZ$, $l\geq 0$) with coefficients in $\bbQ[v]$, let $I(F)$ be the termwise antiderivative of $-F\,\dd x=-F(p^3-vp^{-2})\,\dd p$ with integration constant $0$: the termwise antiderivative of $p^j(\log p)^l\dd p$ is $p^{j+1}$ times a polynomial in $\log p$ if $j\neq -1$, and $(\log p)^{l+1}/(l+1)$ if $j=-1$, so it has no term $p^0(\log p)^0$, and $I(F)$ is the antiderivative without such a term. Put $\phi^b_{k+1}:=I(\phi^b_k)$. Then $D_x\phi^b_{k+1}=-\phi^b_k$, and $\phi^b_k$ is homogeneous of weight $\wt(\phi^b_k)=w_b+4k$, where $w_b:=\nu_b+1$, that is $(w_0,\dots,w_4)=(1,0,4,3,2)$.
\end{defn}

\begin{prop}[The leg dictionary]
\label{prop:mod-LD}
Let $\varrho=(\varrho_0,\dots,\varrho_4):=(-1,1,-1,-3,-2)$. For all $a,b\in\bbZ/5$ and $k,d\geq 0$,
\begin{equation}
\label{eq:mod-LD}
\langle\phi^b_k,\dd\xi_{a,d}\rangle=\sum_c\Res_{c}\phi^b_k\,\dd\xi_{a,d}=-\varrho_b\,\delta_{ab}\,\delta_{dk},
\end{equation}
$c$ running over the branch points and $\log p$ being any branch near each of them. Hence $q_{\phi^b_k}=-\varrho_b\,z^kH^b$, and the TR leg of $v^mz^kH^b$ (the insertion $v^m\psi^kH^b$, \cref{conv:tra-forms}) is $-\varrho_b^{-1}v^m\phi^b_k$. In particular the unit $H^0$ has the TR leg $y=p$; $\Phi_{y\dd x}=p^5/5-v\log p=-\phi^0_1$ has $q_{\Phi_{y\dd x}}=-z\,H^0$; $-\log p$ has $q_{-\log p}=H$; the Euler-dilaton vector $\se=z\,H^0+vH$ has the TR leg $-p^5/5$; and $q_{f(x)}=0$ for $f\in\bbQ(v)[x]$.
\end{prop}

\begin{proof}
\emph{Step 1 (branches and functions of $x$).} Each leg is a polynomial in $\log p$ with coefficients in $\bbQ[v][p,p^{-1}]$, so it is analytic near every branch point $c\neq 0$, for any branch of $\log p$. A change of branch $\log p\mapsto\log p+2\pi\ri N$ changes $\phi^b_k$ by a function $\delta^b_k$ (a local letter) that is a polynomial in $x$ with constant coefficients: $\delta^b_0$ is $0$ or $2\pi\ri N$; the identity $\dd\phi^b_{k+1}=-\phi^b_k\,\dd x$ of \cref{def:mod-legs} is an identity of expressions in $p$ and $\log p$, with $\dd\log p=\dd p/p$, and it survives the substitution, so $\dd\delta^b_{k+1}=-\delta^b_k\,\dd x$; and if $\delta^b_k=P(x)$ with a polynomial $P$, then $\delta^b_{k+1}+\int_0^xP$ has zero differential near $c$, hence is constant. For $f$ holomorphic near $c$, repeated integration by parts gives $\Res_cf(x)\,\dd\xi_{a,d}=\pm\Res_cf^{(d+1)}(x)\,\xi_a\dd x$, which vanishes because $\xi_a\dd x=\varepsilon'_av^{\alpha_a}p^{m_a-2}\dd p$ is holomorphic at $c\neq 0$. So the pairings do not depend on the branches, and $q_{f(x)}=0$.
\emph{Step 2 (lowering).} For $d\geq 1$, integration by parts at each $c$ gives $\langle\phi,\dd\xi_{a,d}\rangle=-\sum_c\Res_c\xi_{a,d}\,\dd\phi=-\langle D_x\phi,\dd\xi_{a,d-1}\rangle$, since $\xi_{a,d}\dd\phi=(D_x\xi_{a,d-1})(D_x\phi)\,\dd x=(D_x\phi)\,\dd\xi_{a,d-1}$. Hence $\langle\phi^b_{k+1},\dd\xi_{a,d}\rangle=\langle\phi^b_k,\dd\xi_{a,d-1}\rangle$.
\emph{Step 3 ($d=0$, $k\geq 1$).} $\langle\phi^b_k,\dd\xi_a\rangle=-\sum_c\Res_c\xi_a\,\dd\phi^b_k=\sum_c\Res_c\phi^b_{k-1}\,\xi_a\dd x=0$, since $\xi_a\dd x$ is holomorphic at $c$.
\emph{Step 4 ($k=0$, $d\geq 1$).} $\langle\phi^b_0,\dd\xi_{a,d}\rangle=-\sum_c\Res_cD_x^d(\xi_a)\,\dd\phi^b_0$. This form is meromorphic on $\bbP^1$ ($\dd\log p=\dd p/p$); it is holomorphic at $p=0$, because $D_x$ maps functions holomorphic at $0$ to $O(p^2)$ while $\dd\phi^b_0=O(p^{-2})\dd p$; and it is $O(p^{1-4d})\dd p$ at $\infty$, because $D_x^d\xi_a=O(p^{-1-4d})$ and $\dd\phi^b_0=O(p^2)\dd p$. By the residue theorem the residues at the branch points add up to $0$.
\emph{Step 5 ($k=d=0$).} Write $\dd\phi^b_0=\gamma_bp^{n_b}\dd p$ with the local constants $(\gamma_b,n_b)=(1,0),(1,-1),(-v,-2),(3,2),(2,1)$. Since $\Res_cp^N\dd p/(p^5-v)=c^{N-4}/5$ and $\sum_cc^N=5v^{N/5}$ if $5\mid N$ and $0$ otherwise,
\[
\langle\phi^b_0,\dd\xi_a\rangle=-\varepsilon'_a\gamma_bv^{\alpha_a}\,\frac15\sum_cc^{m_a+n_b-4},
\]
which is non-zero only if $m_a+n_b\equiv 4\pmod 5$, that is, only if $a=b$. For $a=b=0,\dots,4$ the values are $1,-1,1,3,2$, that is $-\varrho_b$ (for $b=1$: $-v\cdot v^{-1}$; for $b=2$: $-(-v)v^{-1}$).
\emph{Step 6.} If $k=d$, Step 2 applied $k$ times and Step 5 give $-\varrho_b\delta_{ab}$; if $k>d$, Step 2 and Step 3 give $0$; if $k<d$, Step 2 and Step 4 give $0$. This is \eqref{eq:mod-LD}, and $q_{\phi^b_k}=\sum_{a,d}\langle\phi^b_k,\dd\xi_{a,d}\rangle H^az^d=-\varrho_bz^kH^b$. The consequences: $\phi^0_0=p$ and $-\varrho_0=1$; $\phi^0_1=I(p)=-(p^5/5-v\log p)$, so $q_{\Phi_{y\dd x}}=-q_{\phi^0_1}=\varrho_0z\,H^0=-zH^0$; $q_{\log p}=-\varrho_1H=-H$; and the TR leg of $\se$ is $\phi^0_1-v\phi^1_0=-p^5/5$.
\end{proof}

\begin{cor}[Free energies of the model and of the curve]
\label{cor:mod-trF}
For $h\geq 2$ let $\Fm_h(v)=\langle\ \rangle_h(v)=\int_{\overline{\mathcal M}_h}\Omega^{\mathrm{mod}}_{v;h,0}$ be the genus-$h$ free energy of the model CohFT (\cref{def:set-modcohft}), the value at the base point $-z1$ of the dilaton-shifted potential $\overline F_h$ (\cref{def:set-modcohft,def:set-graphsum}). Then $\Fm_h=(-5)^{1-h}\FTR_h$.
\end{cor}

\begin{proof}
On the side of the curve, $(2-2h)\FTR_h=\sum_b\Res_{\varsigma_b}\Phi_{y\dd x}\,\omega_{h,1}=\langle\Phi_{y\dd x}\rangle^{\mathrm{TR}}_h$ by \cref{conv:mod-kernel}; the branch of $\log p$ does not matter, since $\omega_{h,1}$ has zero residues (\cref{lem:mod-xfun}). Only this definition is used on the side of the curve. (It is the definition \cite[Def.~4.3 (4-10)]{Eynard:2007kz} up to the sign $\FTR_h=-F^{(h)}$, and, as remarked after \cite[Thm.~4.7]{Eynard:2007kz}, it extends to $n=0$ the dilaton equation \cite[Thm.~4.7 (4-30)]{Eynard:2007kz}, stated there for $n\geq 1$, which in the convention of \cref{conv:mod-kernel} reads $\sum_b\Res_{\varsigma_b}\Phi_{y\dd x}\,\omega_{h,n+1}(\cdot,J)=(2-2h-n)\,\omega_{h,n}(J)$ for $2h-2+n>0$.) By \cref{thm:mod-tr} and \cref{prop:mod-LD}, $\langle\Phi_{y\dd x}\rangle^{\mathrm{TR}}_h=(-5)^{h-1}\langle q_{\Phi_{y\dd x}}\rangle_h=-(-5)^{h-1}\langle\psi H^0\rangle_h$. The model CohFT has the flat unit $H^0$, and the dilaton equation (\cref{lem:set-dilaton}~(i), (ii), with no further insertion) gives $\langle\psi H^0\rangle_h=(2h-2)\Fm_h$. Comparing, $(2-2h)\FTR_h=(2-2h)(-5)^{h-1}\Fm_h$.
\end{proof}

% Section 6, part 5.
The weight spaces and the modules of legs are needed before the logarithmic part of the legs can be computed; they are used again in \cref{sec:mod-lattice}.

\begin{defn}[Weight spaces, the leg map, the leg class]
\label{def:mod-weights}
For $w\geq 0$ let $V_w$ be the $\bbQ$-span of the insertions $v^m\psi^kH^b$ ($m,k\geq 0$, $b\in\bbZ/5$) with $w_b+4k+5m=w$, and $\mathbb F:=\bigoplus_wV_w$, the free graded $\bbQ[v]$-module on the $\psi^kH^b$ (the weights of \cref{conv:tra-forms}). The \emph{leg map} is the $\bbQ$-linear map $\Lambda\colon\mathbb F\to\{\text{legs}\}$, $\Lambda(v^m\psi^kH^b):=\varrho_b^{-1}v^m\phi^b_k$; by \cref{prop:mod-LD} the TR leg of $X\in\mathbb F$ is $-\Lambda(X)$. Let $\bbQ[p,\bar p]$, $\bar p=v/p$, be the subring of $\bbQ[v][p,p^{-1}]$ spanned by the $v^mp^j$ with $j\geq -m$; recall the leg class $\cL=\bbQ[v][p]+\bbQ(v)[x]$ (\cref{def:mod-pair}). Put
\[
\mathsf M:=\bbQ[p,\bar p]\oplus\log p\cdot\bbQ[v,x],\qquad \mathsf N:=\mathsf M/\bbQ[v,x]=\bbQ[p,\bar p]/\bbQ[v,x]\oplus\log p\cdot\bbQ[v,x],
\]
(the sum is direct since $1$ and $\log p$ are independent over Laurent polynomials; $\bbQ[v,x]\subset\bbQ[p,\bar p]$ as $v=p\bar p$, $x=p^4/4+\bar p$), graded by weight. For $F=A+\log p\cdot B\in\mathsf M$ the element $B\in\bbQ[v,x]$ is the \emph{$\log p$-coefficient} of $F$, and $F$ is \emph{log-free} if $B=0$. Finally $\mathrm d(w):=\lfloor w/5\rfloor+1-[5\mid w]$ and $\mathrm e(w):=\#\{(m,n)\in\bbZ_{\geq 0}^2: 5m+4n=w\}$.
\end{defn}

\begin{lem}[Closure]
\label{lem:mod-closure}
\begin{enumerate}
\item For $a,c\geq 0$, $I(p^a\bar p^c)=-p^{a+4}\bar p^c/(a-c+4)+p^a\bar p^{c+1}/(a-c-1)$, a term with a zero denominator being read as $-v^c\log p$, respectively $+v^{c+1}\log p$.
\item For $g\in\bbQ[v,x]$ with an $x$-primitive $G$, $I(\log p\cdot g(x))=-\log p\cdot G(x)+\int G(x)\,\dd p/p$, where the last integral is termwise: $\int p^a\bar p^c\,\dd p/p=p^a\bar p^c/(a-c)$ if $a\neq c$ and $=v^c\log p$ if $a=c$.
\item $I$ is $\bbQ[v]$-linear, $I(\mathsf M)\subset\mathsf M$ and $I(\bbQ[v,x])\subset\bbQ[v,x]$. Hence every leg $\phi^b_k$ lies in $\mathsf M$, and $\mathsf M$ is a $\bbQ[v,x]$-module.
\end{enumerate}
\end{lem}

\begin{proof}
(i) $-\int v^cp^{a-c}(p^3-vp^{-2})\,\dd p=-v^cp^{a-c+4}/(a-c+4)+v^{c+1}p^{a-c-1}/(a-c-1)$, with $\int\dd p/p=\log p$ in the exceptional cases; and $v^cp^{a-c+4}=p^{a+4}\bar p^c$, $v^{c+1}p^{a-c-1}=p^a\bar p^{c+1}$. (ii) $\dd(\log p\cdot G(x))=G(x)\,\dd p/p+\log p\cdot g(x)\,\dd x$, and $G(x(p))\in\bbQ[p,\bar p]$; the right side has no term $p^0(\log p)^0$ (the termwise $\int G\,\dd p/p$ has none, and $\log p\cdot G$ has no term free of $\log p$), so it is the antiderivative of \cref{def:mod-legs}. A different primitive $G+C$ changes $-\log p\cdot G$ by $-C\log p$ and $\int G\,\dd p/p$ by $C\log p$, so the identity holds for every primitive. (iii) Linearity is clear. By (i), $I$ maps $\bbQ[p,\bar p]$ into $\mathsf M$, the exceptional terms lying in $\log p\cdot\bbQ[v]$. By (ii), $I(\log p\cdot g)\in\mathsf M$, since $G\in\bbQ[v,x]$ and $\int G\,\dd p/p\in\bbQ[p,\bar p]+\log p\cdot\bbQ[v]$. For $g\in\bbQ[v,x]$, $-g(x)\,\dd x=-\dd G(x(p))$ is exact, with no term $p^{-1}\dd p$, so $I(g)=-G(x)+[p^0]G(x(p))\in\bbQ[v,x]$. The five $\phi^b_0$ lie in $\mathsf M$, and induction gives $\phi^b_k\in\mathsf M$. Finally $\bbQ[p,\bar p]$ is a ring containing $\bbQ[v,x]$.
\end{proof}

\begin{lem}[Generating function of the log coefficients]
\label{lem:mod-genfn}
For all $b$ and $k$, the $\log p$-coefficient of $\phi^b_k$ is $\sum_{i=0}^k\mathcal R^b_{k-i}(-x)^i/i!$, where $\mathcal R^b_n:=\Res_{p=0}\Wv^n\,\dd\phi^b_0/n!$, a monomial in $v$.
\end{lem}

\begin{proof}
Fix $b$, write $\phi_k:=\phi^b_k$, and put $\mathcal B_n:=\sum_{j=0}^nx^j\phi_{n-j}/j!$, the coefficient of $\epsilon^n$ in $e^{\epsilon x}\sum_k\epsilon^k\phi_k$ ($\epsilon$ a formal variable). Since $\dd\phi_m/\dd p=-x'\phi_{m-1}$ for $m\geq 1$, the derivative telescopes: $\dd\mathcal B_n/\dd p=x^n\phi_0'/n!$, a Laurent polynomial in $p$. By \cref{lem:mod-closure}~(iii), $\mathcal B_n\in\mathsf M$; write $\mathcal B_n=A+\log p\cdot B$ with $A\in\bbQ[p,\bar p]$ and $B\in\bbQ[v,x]$. Then $\dd\mathcal B_n/\dd p=A'+B/p+\log p\cdot\dd B(x(p))/\dd p$ has no logarithmic term, so $B(x(p))$ does not depend on $p$, and $B\in\bbQ[v]$. As $A'$, a derivative, has no term $p^{-1}$, $B$ is the coefficient of $p^{-1}$ in $x^n\phi_0'/n!$, that is $B=\mathcal R^b_n$. Inverting, $\sum_k\epsilon^k\phi_k=e^{-\epsilon x}\sum_n\epsilon^n\mathcal B_n$, so $\phi_k=\sum_{i=0}^k(-x)^i\mathcal B_{k-i}/i!$; taking $\log p$-coefficients, which commutes with multiplication by $\bbQ[v,x]$, gives the claim. By weight, $\mathcal R^b_n$ is a monomial in $v$.
\end{proof}

\begin{lem}[The residue identity]
\label{lem:mod-resid}
For all $n\geq 0$ and $b$, $\mathcal R^b_n=\varrho_b\,\Pi_0(\psi^nH^b)$.
\end{lem}

\begin{proof}
By \cref{def:gau-van}, $\IC_m=\Res_{p=0}p^me^{-\Wv(p)/z}\dd p$ is a formal residue, taken coefficientwise, so
\[
[z^{-n}]\IC_m=\frac{(-1)^n}{n!}\Res_{p=0}p^m\Wv^n\,\dd p .
\]
With $\YC_b=\varepsilon_b\IC_{\nu_b}$ and $\Pi_0(z^nH^b)=(-1)^{n+1}[z^{-n}]\YC_b$ (\cref{def:tra-tower}) this gives
\[
\Pi_0(\psi^nH^b)=-\frac{\varepsilon_b}{n!}\Res_{p=0}p^{\nu_b}\Wv^n\,\dd p .
\]
Compare with $\dd\phi^b_0=\dd p,\ \dd p/p,\ -vp^{-2}\dd p,\ 3p^2\dd p,\ 2p\,\dd p$: for $b=0,1,3,4$ the identity is immediate, with $\varrho_b=-1,1,-3,-2$. For $b=2$, $\Res_{p=0}\Wv^n(p^3-vp^{-2})\,\dd p=\Res_{p=0}\dd(\Wv^{n+1})/(n+1)=0$, so $\mathcal R^2_n=-v\Res_{p=0}p^{-2}\Wv^n\dd p/n!=-\Res_{p=0}p^3\Wv^n\dd p/n!=-\Pi_0(\psi^nH^2)$, and $\varrho_2=-1$.
\end{proof}

\begin{prop}[Log coefficients of the legs]
\label{prop:mod-dii}
For all $b$ and $k$,
\begin{equation}
\label{eq:mod-dii}
\text{the }\log p\text{-coefficient of }\phi^b_k\ =\ \varrho_b\sum_{i=0}^k\Pi_i(\psi^kH^b)\,\frac{(-x)^i}{i!}.
\end{equation}
Hence, for $X\in V_w$, the $\log p$-coefficient of $\Lambda(X)$ is $\sum_{i\geq 0}\Pi_i(X)(-x)^i/i!$, and $\Lambda(X)$ is log-free if and only if $X\in\sKo$.
\end{prop}

\begin{proof}
By \cref{lem:mod-genfn,lem:mod-resid} and the shift rule $\Pi_i(\psi^kH^b)=\Pi_0(\psi^{k-i}H^b)$ of \cref{def:tra-tower}, the $\log p$-coefficient of $\phi^b_k$ is $\sum_i\varrho_b\Pi_0(\psi^{k-i}H^b)(-x)^i/i!=\varrho_b\sum_i\Pi_i(\psi^kH^b)(-x)^i/i!$. For $X=\sum_ta_tv^{m_t}\psi^{k_t}H^{b_t}\in V_w$ the factors $\varrho_b^{-1}$ of $\Lambda$ cancel, and since the $\Pi_i$ are $\bbQ[v]$-linear the $\log p$-coefficient of $\Lambda(X)$ is $\sum_i\Pi_i(X)(-x)^i/i!$. The $\Pi_i(X)$ lie in $\bbQ[v]$ and the powers of $x$ are linearly independent over $\bbQ(v)$, so $\Lambda(X)$ is log-free if and only if $\Pi_i(X)=0$ for all $i\geq 0$, that is (\cref{def:tra-slice}) if and only if $X\in\sKo$.
\end{proof}

The normalising constants $\varrho_b$ in $\Lambda$ matter for $\sKo$ as a subspace: $\se=z\,H^0+vH$ lies in $\sKo$ and $\Lambda(\se)=p^5/5$ is log-free, whereas $\phi^0_1+v\phi^1_0=-p^5/5+2v\log p$ is not.

% Section 6, part 6.
\subsection{The lattice theorem}
\label{sec:mod-lattice}

\begin{lem}[Counting]
\label{lem:mod-count}
For every $w\geq 0$, $V_w$ has exactly $\mathrm d(w)$ basis triples $(m,k,b)$ with $b\neq 1$ and exactly $\mathrm e(w)$ with $b=1$.
\end{lem}

\begin{proof}
Fix $m$ with $5m\leq w$ and put $r:=w-5m$. A triple with $b\neq 1$ needs $w_b+4k=r$ with $w_b\in\{1,4,3,2\}$ for $b=0,2,3,4$. These four weights meet every class modulo $4$ exactly once, so there is exactly one candidate $b$, and it works unless $w_b>r$; this happens only for the class $0$, where $w_2=4$ and $r=0$, that is $w=5m$. Hence the number is $\lfloor w/5\rfloor+1-[5\mid w]=\mathrm d(w)$. The triples with $b=1$ ($w_1=0$) are the pairs $(m,k)$ with $5m+4k=w$.
\end{proof}

\begin{lem}[The leg class]
\label{lem:mod-classL}
\begin{enumerate}
\item $\bbQ[p,\bar p]=\bbQ[p,x]=\bbQ[v][p]+\bbQ[v][x]$, and $\bbQ[p,x]$ is a free $\bbQ[v,x]$-module with basis $1,p,p^2,p^3,p^4$.
\item $\bbQ[v][p^{\pm1}]\cap\bbQ(v)[x]=\bbQ[v][x]$.
\item $\cL\cap\bbQ[v][p^{\pm1}]=\bbQ[p,\bar p]$: a Laurent polynomial with coefficients in $\bbQ[v]$ lies in $\cL$ if and only if each of its terms $v^mp^j$ has $j\geq -m$.
\item $\cL/\bbQ(v)[x]=\bbQ[p,\bar p]/\bbQ[v,x]=\bbQ[v][p]/\bbQ[v]=\bigoplus_{i=1}^4\bbQ[v,x]\,p^i$; it is free over $\bbQ[v]$, of dimension $\mathrm d(w)$ in weight $w$.
\end{enumerate}
\end{lem}

\begin{proof}
(i) $\bar p=x-p^4/4$ gives $\bbQ[p,\bar p]=\bbQ[p,x]$, and $v=p\bar p=px-p^5/4$. The elements $p$ and $x$ are algebraically independent, since $p$ and $v$ are and $v=px-p^5/4$. The map $\bbQ[v,x][T]/(T^5-4xT+4v)\to\bbQ[p,x]$, $T\mapsto p$, is onto; its source is $\bbQ[x,T]$ (eliminate $v=(4xT-T^5)/4$), so by the independence it is an isomorphism, which gives the basis. Next, $\bbQ[v][p]+\bbQ[v][x]\subset\bbQ[p,\bar p]$. Conversely let $F\in\bbQ[p,\bar p]$ be homogeneous with a negative power of $p$, and let $v^mp^{-n}$ ($n>0$, $m\geq n$) be its term with the most negative exponent; by homogeneity it is the only term with this exponent. Since $x^n=\sum_t\binom nt4^{t-n}v^tp^{4n-5t}$ has the lowest term $v^np^{-n}$, subtracting a multiple of $v^{m-n}x^n$ gives an element of $\bbQ[p,\bar p]$ whose most negative exponent is larger; by induction $F\in\bbQ[v][p]+\bbQ[v][x]$. (ii) Let $P=\sum_{n\leq N}a_nx^n$ with $a_N\neq 0$ lie in $\bbQ[v][p^{\pm1}]$. Only $x^N$ contributes to $p^{4N}$, with the coefficient $a_N4^{-N}$, so $a_N\in\bbQ[v]$; induct on $N$. (iii) The inclusion $\supset$ follows from (i). If $F=R+P$ with $R\in\bbQ[v][p]$ and $P\in\bbQ(v)[x]$ is a Laurent polynomial over $\bbQ[v]$, then $P=F-R\in\bbQ[v][x]$ by (ii), so $F\in\bbQ[p,\bar p]$ by (i). (iv) The map $\bbQ[p,\bar p]\to\cL/\bbQ(v)[x]$ is onto because $\bbQ[v][p]\subset\bbQ[p,\bar p]$, with kernel $\bbQ[p,\bar p]\cap\bbQ(v)[x]=\bbQ[v,x]$ by (ii). By (i), $\bbQ[p,\bar p]/\bbQ[v,x]=\bigoplus_{i=1}^4\bbQ[v,x]p^i$. As the lowest term of $x^n$ is $v^np^{-n}$, $\bbQ[v][p]\cap\bbQ[v][x]=\bbQ[v]$, so the quotient is also $\bbQ[v][p]/\bbQ[v]$, whose weight-$w$ part has the basis $v^mp^{w-5m}$, $0\leq m\leq w/5$, $5m<w$: $\mathrm d(w)$ elements.
\end{proof}

\begin{lem}[Reduction modulo $v$]
\label{lem:mod-modv}
In every weight, $\Lambda\otimes\bbQ[v]/(v)\colon\mathbb F/v\mathbb F\to\mathsf N/v\mathsf N$ is an isomorphism; here $\Lambda$ is followed by the projection $\mathsf M\to\mathsf N$.
\end{lem}

\begin{proof}
$\Lambda$ is graded and $\bbQ[v]$-linear with values in $\mathsf M$ (\cref{lem:mod-closure}). \emph{Source.} In weight $w$, $\mathbb F/v\mathbb F$ has the basis $\psi^kH^b$ with $w_b+4k=w$: one element if $4\nmid w$; $\{\psi^{n-1}H^2,\psi^nH\}$ if $w=4n\geq 4$; $\{H\}$ if $w=0$. \emph{Target.} $\mathsf N/v\mathsf N=\bbQ[p,\bar p]/(\bbQ[v,x]+p\bar p\,\bbQ[p,\bar p])\oplus\log p\cdot\bbQ[x]$. The ring $\bbQ[p,\bar p]$ is the polynomial ring in $p,\bar p$, and $\bbQ[p,\bar p]/(p\bar p)$ has the basis $1$, $p^a$, $\bar p^c$ ($a,c\geq 1$). The image of $\bbQ[v,x]$ is spanned by $1$ and the $x^n\equiv p^{4n}/4^n+\bar p^n$; so the first summand has the basis $\{p^a:a\geq 1\}$, with $\bar p^n\equiv -4^{-n}p^{4n}$, and $(\mathsf N/v\mathsf N)_w$ has the basis $\{p^w\}$ if $w\geq 1$, $4\nmid w$; $\{p^{4n},\log p\cdot x^n\}$ if $w=4n\geq 4$; $\{\log p\}$ if $w=0$. \emph{The map.} By \cref{lem:mod-closure}~(iii), $I$ induces a map on $\mathsf N/v\mathsf N$. By \cref{lem:mod-closure}~(i), $I(p^a)\equiv -p^{a+4}/(a+4)$ modulo $v\mathsf M$ for $a\geq 1$, so $\phi^b_k\equiv(-1)^kp^{w_b+4k}/\prod_{i=1}^k(w_b+4i)$ for $b=0,4,3$; since $\bar p\equiv -p^4/4$ modulo $\bbQ[v,x]$, $\phi^2_k\equiv(-1)^{k+1}p^{4k+4}/\big(4\prod_{i=1}^k(4+4i)\big)$; and by \cref{lem:mod-closure}~(ii) the $\log p$-component of $\phi^1_k$ is $(-1)^kx^k/k!$ modulo $v$. In weight $4n$ the matrix in these bases is triangular with non-zero diagonal entries; in the other weights it is a non-zero $1\times 1$ matrix (the factors $\varrho_b^{-1}$ do not matter).
\end{proof}

\begin{thm}[The lattice theorem]
\label{thm:mod-lattice}
For every $w\geq 0$:
\begin{enumerate}
\item[(a)] $V_w$ has $\mathrm d(w)$ basis triples with $b\neq 1$ and $\mathrm e(w)$ with $b=1$;
\item[(b)] for $X\in\sKo\cap V_w$, $\Lambda(X)$ lies in $\bbQ[p,\bar p]=\bbQ[v][p]+\bbQ[v][x]$, exactly and not only modulo functions of $x$;
\item[(c)] $\Lambda\colon\sKo\cap V_w\to(\cL/\bbQ(v)[x])_w$ is bijective, and both sides have dimension $\mathrm d(w)$;
\item[(d)] $\Lambda\colon V_w\to\mathsf N_w$ is bijective; $\mathsf N$ is a free $\bbQ[v]$-module with basis $\{\Lambda(\psi^kH^b)\}$.
\end{enumerate}
The coefficients must lie in $\bbQ[v]$: over $\bbQ(v)$ the statement fails, since $\Lambda(H^2/v)=-1/p\notin\cL$.
\end{thm}

\begin{proof}
(a) is \cref{lem:mod-count}. (b) $\Lambda(X)\in\mathsf M$ by \cref{lem:mod-closure}, and it is log-free by \cref{prop:mod-dii}, so $\Lambda(X)\in\bbQ[p,\bar p]$, which equals $\bbQ[v][p]+\bbQ[v][x]$ by \cref{lem:mod-classL}~(i). (d) $\Lambda\colon\mathbb F\to\mathsf N$ is a graded $\bbQ[v]$-linear map between modules with finite-dimensional homogeneous pieces, zero in negative weights. Let $C$ be its cokernel. By right exactness of $\otimes\bbQ[v]/(v)$ and \cref{lem:mod-modv}, $C/vC=0$. If $C\neq 0$, let $w$ be minimal with $C_w\neq 0$; then $(vC)_w=vC_{w-5}=0$, so $C_w=(C/vC)_w=0$, a contradiction. So $\Lambda$ is onto. Since $\dim V_w=\mathrm d(w)+\mathrm e(w)$ (\cref{lem:mod-count}) and $\dim\mathsf N_w=\mathrm d(w)+\mathrm e(w)$ (\cref{lem:mod-classL}~(iv), and $\dim\bbQ[v,x]_w=\mathrm e(w)$ as $v,x$ are algebraically independent), $\Lambda_w$ is bijective and $\mathsf N\cong\mathbb F$ is free. (c) By (b) the image lies in $\bbQ[p,\bar p]$, and $\cL/\bbQ(v)[x]=\bbQ[p,\bar p]/\bbQ[v,x]$ by \cref{lem:mod-classL}~(iv). If $\Lambda(X)\in\bbQ(v)[x]$, then $\Lambda(X)\in\bbQ[p,\bar p]\cap\bbQ(v)[x]=\bbQ[v,x]$ (\cref{lem:mod-classL}~(ii)), so $\Lambda(X)=0$ in $\mathsf N$ and $X=0$ by (d). Given $\phi\in(\bbQ[p,\bar p]/\bbQ[v,x])_w$, the element $X:=\Lambda^{-1}(\phi\oplus 0)$ of (d) has a log-free leg, so $X\in\sKo$ by \cref{prop:mod-dii}. Both sides have dimension $\mathrm d(w)$ by \cref{lem:mod-classL}~(iv). Finally $\Lambda(H^2)=\varrho_2^{-1}v/p=-v/p$, and $-1/p\notin\cL$ by \cref{lem:mod-classL}~(iii).
\end{proof}

\begin{cor}[Splitting of $V_w$]
\label{cor:mod-split}
$V_w=(\sKo\cap V_w)\oplus\operatorname{span}\{v^m\psi^kH:5m+4k=w\}$. Every $X\in V_w$ is $X_0+\sum a_{m,k}v^m\psi^kH$ with $X_0\in\sKo$, the $a_{m,k}$ being read off from the $\log p$-coefficient of $\Lambda(X)$.
\end{cor}

\begin{proof}
The decomposition is \cref{lem:tra-frames}~(a) in weight $w$. Alternatively: by \cref{lem:mod-closure}~(ii) and induction, the $\log p$-coefficient of $\phi^1_k$ is $(-1)^kx^k/k!$ plus terms of lower degree in $x$, so the map that sends $X$ to the $\log p$-coefficient of $\Lambda(X)$, restricted to the span of the $v^m\psi^kH$, is triangular onto $\bbQ[v,x]_w$ ($v^m\psi^kH\leftrightarrow v^mx^k$), hence bijective; its kernel on $V_w$ is $\sKo\cap V_w$ by \cref{prop:mod-dii}.
\end{proof}

\begin{cor}[Splitting of $\sK\cap V_w$]
\label{cor:mod-Ksplit}
$\sK\cap V_w=(\sKo\cap V_w)\oplus\bbQ v^{w/5}H$ if $5\mid w$, and $\sK\cap V_w=\sKo\cap V_w$ otherwise. On $\sK$ the only legs beyond $\bbQ[p,\bar p]$ are multiples of $\log p$: for $X\in\sK\cap V_w$, $\Lambda(X)=\Lambda(X-\Pi_0(X)H)+\Pi_0(X)\log p$ with $\Lambda(X-\Pi_0(X)H)\in\bbQ[p,\bar p]$.
\end{cor}

\begin{proof}
By \cref{prop:mod-dii}, $X\in\sK$ if and only if the $\log p$-coefficient $\sum_i\Pi_i(X)(-x)^i/i!$ of $\Lambda(X)$ does not involve $x$, that is, equals $\Pi_0(X)$, which lies in $\bbQ v^{w/5}$ (\cref{def:tra-tower}; zero unless $5\mid w$). Write $X=X_0+Y$ as in \cref{cor:mod-split}. Then $\Lambda(Y)$ has the same $\log p$-coefficient as $\Lambda(X)$, while $\Lambda(\Pi_0(X)H)=\Pi_0(X)\log p$ has the $\log p$-coefficient $\Pi_0(X)$; by injectivity on the span of the $v^m\psi^kH$, $Y=\Pi_0(X)H$. Conversely $\sKo\subset\sK$ and $H\in\sK$. The last statement follows from \cref{thm:mod-lattice}~(b), as $X-\Pi_0(X)H\in\sKo\cap V_w$.
\end{proof}

\begin{cor}[A basis of $\sKo\cap V_w$]
\label{cor:mod-basis}
The elements $v^m\fr(k,a)$ ($a\neq 1$) of weight $w$ form a basis of $\sKo\cap V_w$, and $\dim\sKo\cap V_w=\mathrm d(w)$.
\end{cor}

\begin{proof}
By \cref{lem:tra-frames}~(b), (c) in the graded setting, the frame coordinates identify $\sKo\cap V_w$ with the families $(c_{(k,a)})_{a\neq 1}$ in which $c_{(k,a)}$ is a rational multiple of $v^{(w-4k-\nu_a-1)/5}$, zero unless the exponent is a non-negative integer; so the $v^m\fr(k,a)$ of weight $w$, that is the triples $(m,k,a)$ with $a\neq 1$ and $w_a+4k+5m=w$, form a basis. Their number is $\mathrm d(w)$ by \cref{thm:mod-lattice}~(a). (The same dimension follows from \cref{thm:mod-lattice}~(c).)
\end{proof}

% Section 6, part 7.
\subsection{L-regularity and the Gap Theorem}
\label{sec:mod-Lreg}

We use the closed formula of \cite[Thm.~7.1]{ABDKS25}: for an arbitrary rational function $x(z)$ on $\bbP^1$ such that $\dd x$ has only simple zeros, the differentials of topological recursion for the curve $(x=x(z),y=z)$ with the Bergman kernel and the initial form $\omega^{(0)}_{1,0}=-y\,\dd x$ are
\begin{equation}
\label{eq:mod-abdks}
\omega^{(g)}_{m,0}=[\hbar^{2g}]\sum_{r_1,\dots,r_m\geq 0}\prod_{i=1}^m\big(\mathcal D_{x_i}\big)^{r_i}[w_i^{r_i}]\sum_{\Gamma}\hbar^{2g(\Gamma)}\prod_{\{i,j\}\in E(\Gamma)}W_{i,j}\prod_{i=1}^mW_i
\end{equation}
\cite[(7.9)]{ABDKS25}, where $\Gamma$ runs over the connected simple graphs on the vertices $1,\dots,m$, $g(\Gamma)$ is the first Betti number, $\mathcal D_{x_i}\eta:=\dd(\eta/\dd x_i)$ on $1$-forms in $z_i$, $S(u):=\sinh(u/2)/(u/2)$,
\begin{gather*}
W_i:=-\exp\!\big(-w_i(S(w_i\hbar\partial_{z_i})-1)\,x(z_i)\big)\frac{\dd z_i}{w_i},\\
W_{i,j}:=\frac{1}{\hbar^2}\Big(\exp\!\big(\hbar^2w_iw_jS(w_i\hbar\partial_{z_i})S(w_j\hbar\partial_{z_j})(z_i-z_j)^{-2}\big)-1\Big)
\end{gather*}
\cite[(7.6), (7.7)]{ABDKS25}, and $W_{i,j}=\sum_{k\geq 0}(\hbar^2/4)^k(w_i+w_j)^{2k}w_iw_j(z_i-z_j)^{-2k-2}$; for one point \cite[(7.10)]{ABDKS25},
\begin{equation}
\label{eq:mod-710}
\sum_{g\geq 1}\hbar^{2g}\omega^{(g)}_{1,0}=\sum_{r\geq 0}\mathcal D_x^r[w^r]W_1 .
\end{equation}
The standing assumptions of \cite[\S1.1]{ABDKS25} (a compact curve, $x,y$ meromorphic, the critical points of $x$ simple, $y$ regular with $\dd y\neq 0$ there) hold for $x=\Wv$, $y=p$, $v\neq 0$: $\Wv$ is rational, $\dd\Wv=(p^5-v)\dd p/p^2$ has exactly five zeros, all simple, and $y=p$ has $\dd y=\dd p\neq 0$. The kernel of \cite[(1.1)]{ABDKS25} with $\omega^{(0)}_{1,0}=-y\,\dd x$ is the negative of ours (\cref{conv:mod-kernel}), so $\omega_{h,n}=(-1)^n\omega^{(h)}_{n,0}$.

\begin{lem}[Structure of the correlators]
\label{lem:mod-abdks}
For $2h-2+n>0$ there are an integer $K\geq 0$ and $N_{h,n}\in\bbQ[v][p_1,\dots,p_n]$ with
\[
\omega_{h,n}=N_{h,n}(p;v)\,\frac{\prod_i\dd p_i}{\prod_i(p_i^5-v)^K}.
\]
\end{lem}

\begin{proof}
(1) \emph{Each term of \eqref{eq:mod-abdks} lies in $\bbQ[v][p_i,p_i^{-1},(p_i^5-v)^{-1},(p_i-p_j)^{-1}]\prod_i\dd p_i$, and for fixed $(h,n)$ there are finitely many terms.} Since $S(u)-1=\sum_{k\geq 1}u^{2k}/(4^k(2k+1)!)$,
\[
-w\big(S(w\hbar\partial_p)-1\big)\Wv(p)=\sum_{k\geq 1}w^{2k+1}\hbar^{2k}\epsilon_k(p),\qquad \epsilon_k:=-\frac{\Wv^{(2k)}}{4^k(2k+1)!}\quad(\text{local letters}),
\]
and $\Wv^{(2k)}=(p^4/4)^{(2k)}+(2k)!\,v\,p^{-2k-1}\in\bbQ[v][p,p^{-1}]$. In $W_i$ every $\hbar^{2k}$ comes with $w_i^{2k+1}$, and in $W_{i,j}$ every $\hbar^2$ comes with $(w_i+w_j)^2$ and one more factor $(p_i-p_j)^{-2}$ ($z_i=p_i$). So at fixed order $\hbar^{2h}$ the degrees in the $w_i$ are bounded, only finitely many $r_i$ and graphs contribute, and the coefficients are polynomials over $\bbQ$ in the $\epsilon_k(p_i)$ and the $(p_i-p_j)^{-1}$. Finally $\mathcal D_x(f\,\dd p)=\dd\big(p^2f/(p^5-v)\big)$ keeps this ring and only adds powers of $(p^5-v)^{-1}$. (2) \emph{Clearing.} Write the finite sum as $\omega_{h,n}=P\prod_i\dd p_i\big/\big(\prod_ip_i^{a}\prod_i(p_i^5-v)^{K}\prod_{i<j}(p_i-p_j)^{e}\big)$ with $P\in\bbQ[v][p]$. By \cref{conv:mod-kernel}, $\omega_{h,n}$ has poles only at the branch points; $p_i=0$ and the diagonals contain none ($v\neq 0$). So for every $v\neq 0$ the polynomial $P$ vanishes to order at least $e$ along $p_i=p_j$ and at least $a$ along $p_i=0$. Divide $P$ by $(p_i-p_j)^e$, respectively $p_i^a$, as polynomials in $p_i$ over $\bbQ[v][p_k:k\neq i]$ (the divisors are monic, so the division algorithm works over this ring): the remainder vanishes for every $v\neq 0$ and every $p$, hence it is $0$. One factor at a time, all these denominators cancel in $\bbQ[v][p]$, and no negative power of $v$ arises.
\end{proof}

\begin{thm}[L-regularity]
\label{thm:mod-Lreg}
Let $2h-2+n>0$. For legs $\phi_1,\dots,\phi_n\in\bbQ[v][p]$ the residue pairing $\langle\phi_1\cdots\phi_n\rangle^{\mathrm{TR}}_h$ lies in $\bbQ[v]$. For legs in $\cL=\bbQ[v][p]+\bbQ(v)[x]$ the residue pairing lies in $\bbQ[v]$; the components in $\bbQ(v)[x]$ pair to zero.
\end{thm}

\begin{proof}
(1) Fix all variables but $p_l$ off the branch points. By \cref{lem:mod-abdks}, $\phi_l(p_l)\omega_{h,n}$ is a rational $1$-form in $p_l$ with poles at the branch points and possibly at $p_l=\infty$; it is holomorphic at $p_l=0$. By the residue theorem, $\sum_b\Res_{p_l=\varsigma_b}=-\Res_{p_l=\infty}$. (2) The result is again a polynomial over $\prod_{m\neq l}(p_m^5-v)^K$: in one variable, $(p^5-v)^{-K}=\sum_{j\geq 0}\binom{K+j-1}{j}v^jp^{-5K-5j}$ for $|p|^5>|v|$, so
\[
\Res_{p=\infty}\frac{p^a\,\dd p}{(p^5-v)^K}=\begin{cases}-\binom{K+j-1}{j}v^j&\text{if }a+1=5(K+j)\text{ with an integer }j\geq 0,\\ 0&\text{otherwise,}\end{cases}
\]
an element of $\bbQ[v]$. The iterated residues commute, as there are no diagonal poles. So $\langle\phi_1\cdots\phi_n\rangle^{\mathrm{TR}}_h=(-1)^n\Res_{p_1=\infty}\cdots\Res_{p_n=\infty}\big[\prod_l\phi_l(p_l)\,N_{h,n}\prod_l\dd p_l/\prod_l(p_l^5-v)^K\big]$ is a finite $\bbQ[v]$-combination of such residues, hence in $\bbQ[v]$. (3) By \cref{lem:mod-xfun} the components in $\bbQ(v)[x]$ of legs in $\cL$ pair to zero, whatever decomposition is chosen.
\end{proof}

By \cref{lem:mod-scale}~(ii) a pairing with monomial legs $p^{m_l}$ is $Cv^e$ with $e=2-2h-n+\sum_lm_l/5$; \cref{thm:mod-Lreg} says that $C=0$ whenever $e<0$.

\begin{prop}[The period of $\omega_{h,1}$]
\label{prop:mod-period}
Let $h\geq 1$. Then:
\begin{enumerate}
\item $\omega_{h,1}=\dd\Phi_h$ with $\Phi_h$ rational and holomorphic at $p=0$ and $p=\infty$;
\item $\int_\infty^0\omega_{h,1}=\frac{B_{2h}}{2h}\,v^{1-2h}$ along any path avoiding the branch points;
\item for $h\geq 2$, $\FTR_h=\dfrac{B_{2h}}{2h(2h-2)}\,v^{2-2h}$.
\end{enumerate}
\end{prop}

\begin{proof}
\emph{Step 1 (exactness).} Write $\omega^{\mathrm A}_{h,1}:=\omega^{(h)}_{1,0}=-\omega_{h,1}$. By \eqref{eq:mod-710}, $\omega^{\mathrm A}_{h,1}=[\hbar^{2h}]\sum_{r\geq 1}\mathcal D_x^r[w^r]W_1$; the term $r=0$ is $-[w^1]e^{E}\dd p=0$, where $E:=\sum_kw^{2k+1}\hbar^{2k}\epsilon_k=O(w^3)$. For $r\geq 1$, $\mathcal D_x^r(a\,\dd p)=\dd\big(D_x^{r-1}(a/x')\big)$ with $D_x=(1/x')\,\dd/\dd p$ on functions (\cref{def:mod-xi}). Hence $\omega^{\mathrm A}_{h,1}=\dd\Phi^{\mathrm A}_h$ with $\Phi^{\mathrm A}_h:=\sum_{r\geq 1}D_x^{r-1}(a_r/x')$, $a_r:=-[\hbar^{2h}][w^{r+1}]e^{E}=-\sum\frac1{j!}\prod_{i=1}^j\epsilon_{k_i}$, the sum over $j\geq 1$ and $k_1+\dots+k_j=h$ with $r=2h-1+j$; it is finite ($1\leq j\leq h$). Put $\Phi_h:=-\Phi^{\mathrm A}_h$.

\emph{Step 2 ($p=\infty$).} $\deg\epsilon_k\leq 4-2k$, $1/x'=p^2/(p^5-v)$ has degree $-3$, and $D_x$ lowers the degree by $4$; every term of $\Phi^{\mathrm A}_h$ has degree at most $(4j-2h)-3-4(2h-2+j)=5-10h\leq -5$. So $\Phi^{\mathrm A}_h$ is holomorphic at $\infty$ and $\Phi^{\mathrm A}_h(\infty)=0$.

\emph{Step 3 ($p=0$).} Write $\epsilon_k=\epsilon^0_k+(\text{terms of order}\geq 0)$ with $\epsilon^0_k:=-v\,p^{-2k-1}/(4^k(2k+1))$; further $1/x'=-(p^2/v)(1-p^5/v)^{-1}$, and $D_x=\delta_0+(\text{an operator raising the order by at least }6)$ with $\delta_0:=-(p^2/v)\,\dd/\dd p$. In the \emph{pure part} of a term, all $\epsilon_k$, $1/x'$, $D_x$ are replaced by $\epsilon^0_k$, $-p^2/v$, $\delta_0$. In the variable $\bar p=v/p$ one has $\delta_0=\dd/\dd\bar p$, $-p^2/v=-v/\bar p^2$ and $\epsilon^0_k=-\bar p^{2k+1}v^{-2k}/(4^k(2k+1))$. With the local letters $\tilde w:=w\bar p$, $\hbar':=\hbar/(2v)$, $E_0:=-\sum_{k\geq 1}\hbar'^{2k}\tilde w^{2k+1}/(2k+1)$ and $C_{h,r}:=[\hbar'^{2h}][\tilde w^{r+1}]\exp(E_0)$, the pure part of $a_r$ is $-(2v)^{-2h}\bar p^{r+1}C_{h,r}$, and the pure part of the $r$-th term is $(\dd/\dd\bar p)^{r-1}\big[(2v)^{-2h}v\,C_{h,r}\bar p^{r-1}\big]=v(2v)^{-2h}(r-1)!\,C_{h,r}$, a constant. Replacing a pure factor by its correction raises the order in $p$ by at least $3$ (by at least $2k+1$ for $\epsilon_k$, at least $5$ for $1/x'$ and for $D_x$), and $\delta_0$ never lowers the order. So every term is a constant plus $O(p^3)$: $\Phi^{\mathrm A}_h$ is holomorphic at $0$, and $\Phi^{\mathrm A}_h(0)=v(2v)^{-2h}\Sigma_h$ with $\Sigma_h:=\sum_r(r-1)!\,C_{h,r}$.

\emph{Step 4 (the identity $\Sigma_h=-4^hB_{2h}/(2h)$).} Let $\mathsf a(\mathsf x):=\sum_{k\geq 1}\mathsf x^k/(2k+1)$, so that $E_0=-\tilde w\,\mathsf a(\hbar'^2\tilde w^2)$. With $r+1=2h+j$, the coefficient $C_{h,r}$ comes from $E_0^j/j!$ alone, and $C_{h,r}=\frac{(-1)^j}{j!}[\mathsf x^h]\mathsf a(\mathsf x)^j$. Hence, by the binomial series $(1+\mathsf a)^{1-2h}=\sum_j(-1)^j\binom{2h-2+j}{j}\mathsf a^j$ (the term $j=0$ has no $\mathsf x^h$, as $h\geq 1$, and the terms $j>h$ none either),
\[
\Sigma_h=\sum_{j=1}^h(-1)^j\frac{(2h-2+j)!}{j!}[\mathsf x^h]\mathsf a^j=(2h-2)!\,[\mathsf x^h](1+\mathsf a(\mathsf x))^{1-2h}.
\]
With $\mathsf x=\mathsf y^2$, $1+\mathsf a(\mathsf y^2)=\operatorname{artanh}(\mathsf y)/\mathsf y$, so $[\mathsf x^h](1+\mathsf a)^{1-2h}=[\mathsf y^{1}]\operatorname{artanh}(\mathsf y)^{1-2h}$. By the Lagrange--B\"urmann formula for $\mathsf y=\tanh\mathsf u$, $[\mathsf y^n]H(\mathsf u(\mathsf y))=\frac1n\Res_{\mathsf u=0}H'(\mathsf u)\tanh(\mathsf u)^{-n}\dd\mathsf u$ for $n\neq 0$ and every Laurent series $H$ (change of variables and one integration by parts), applied to $H(\mathsf u)=\mathsf u^{1-2h}$ and $n=1$: $[\mathsf y^1]\mathsf u^{1-2h}=(1-2h)[\mathsf u^{2h}](\mathsf u\coth\mathsf u)=(1-2h)4^hB_{2h}/(2h)!$. So $\Sigma_h=-4^hB_{2h}/(2h)$ and $\Phi^{\mathrm A}_h(0)=-\frac{B_{2h}}{2h}v^{1-2h}$. This proves (i).

\emph{Step 5.} (ii) Since $\Phi_h$ is rational, $\int_\infty^0\omega_{h,1}=\Phi_h(0)-\Phi_h(\infty)=-\Phi^{\mathrm A}_h(0)=\frac{B_{2h}}{2h}v^{1-2h}$ along any path avoiding the branch points. (iii) $\Phi_{y\dd x}=p^5/5-v\log p$, so $(2-2h)\FTR_h=\frac15\langle p^5\rangle^{\mathrm{TR}}_h+v\langle-\log p\rangle^{\mathrm{TR}}_h$, the residue pairings with any branch of $\log p$ (\cref{def:mod-pair}). Here $\langle p^5\rangle^{\mathrm{TR}}_h\in\bbQ[v]$ by \cref{thm:mod-Lreg}, of exponent $2-2h<0$ (\cref{lem:mod-scale}~(ii)), hence $0$; and $\langle-\log p\rangle^{\mathrm{TR}}_h=\int_{\mathfrak r}\omega_{h,1}=-\int_\infty^0\omega_{h,1}=-\frac{B_{2h}}{2h}v^{1-2h}$ by (ii). So $\FTR_h=\frac{-v}{2-2h}\cdot\frac{B_{2h}}{2h}v^{1-2h}=\frac{B_{2h}}{2h(2h-2)}v^{2-2h}$.
\end{proof}

\begin{lem}[The vanishing period]
\label{lem:mod-vanper}
For $h\geq 1$, $\langle-\log p\rangle^{\mathrm{TR}}_h=\int_{\mathfrak r}\omega_{h,1}=\sum_b\Res_{\varsigma_b}(-\log p)\,\omega_{h,1}=-\frac{B_{2h}}{2h}\,v^{1-2h}$.
\end{lem}

\begin{proof}
The first equality is \cref{def:mod-pair}. For the second, $\log p$ is single-valued near each $\varsigma_b$ and $\Res_{\varsigma_b}(\log p\,\dd\Phi_h)=-\Res_{\varsigma_b}\Phi_h\,\dd p/p$; the residue theorem for the rational form $\Phi_h\,\dd p/p$, whose other poles are $0$ and $\infty$, gives $\sum_b\Res_{\varsigma_b}\log p\cdot\omega_{h,1}=\Phi_h(0)-\Phi_h(\infty)=\int_\infty^0\omega_{h,1}$, while $\int_{\mathfrak r}\omega_{h,1}=-\int_\infty^0\omega_{h,1}$ for the ray $\mathfrak r$ from $0$ to $\infty$. The value is \cref{prop:mod-period}~(ii).
\end{proof}

\begin{thm}[The Gap Theorem]
\label{thm:mod-gap}
On the model curve, for every $h\geq 2$, $\FTR_h=\dfrac{B_{2h}}{2h(2h-2)}v^{2-2h}$ in the convention of \cref{conv:mod-kernel}; hence $\Fm_h=c_hv^{2-2h}$ with $c_h=(-5)^{1-h}\dfrac{B_{2h}}{2h(2h-2)}=\dfrac{(-1)^{h-1}B_{2h}}{2h(2h-2)\,5^{h-1}}$.
\end{thm}

\begin{proof}
\Cref{prop:mod-period}~(iii) and \cref{cor:mod-trF}.
\end{proof}

\begin{rmk}[A first proof through a formula of Hock]
\label{rmk:mod-hock}
The value of $\FTR_h$ can also be read off from \cite[Cor.~1.2, (1.3)]{Hock25}, whose kernel and free energy \cite[(2.1), (1.1)]{Hock25} are those of \cref{conv:mod-kernel}; for $y$ unramified (here $y=p$ is a global coordinate) it expresses $\FTR_h$ as $\frac{1}{2-2h}\sum\Res\,[\hbar^{2h-1}]$ of an explicit series, the residues being taken at the logarithmic singularities of $x$ (there are none here) and at the other singularities of $x$ and $y$, here the poles $p=0$ and $p=\infty$. The residue at $\infty$ vanishes by a degree count, and at $p=0$ only the leading terms contribute, giving $v^{2-2h}A_h$ with $A_h=(2h-2)!\,[\mathsf x^h](1+\mathsf a(\mathsf x/4))^{1-2h}=4^{-h}\Sigma_h=-B_{2h}/(2h)$ by the Lagrange inversion of Step~4 above, so that $\FTR_h=A_hv^{2-2h}/(2-2h)$. This argument is only sketched here and is not used. \draftnote{The closed forms printed in \cite[\S3.2]{Hock25} are not used; only Cor.~1.2 there is cited.}
\end{rmk}

\input{sections/s06_model_8}



